\subsection{启动与根文件系统}\label{sec:boot}

机器人上的所有软件都来自一张 SD 卡。这块板子上电后，固化在芯片里的一小段引导程序先把
U-Boot 从卡上装入并运行，U-Boot 再把内核镜像读进内存并交出控制权。内核拿到控制权后做的头
等大事，是把 SD 卡上那块以 ext4 格式存放的根文件系统挂载起来，应用程序、模型权重与配置才
有处可寻。Starry 在这里与 Linux 走同一条路，它要挂载的正是 Linux 也在用的那一块 ext4 根分
区，存储介质、分区布局与文件系统格式都不变，变的只是挂载它的内核。

SD 卡这类块设备并不能像内存那样随取随到，向它发起一次读取要经过排队、提交、等待硬件搬运、
最后收取完成信号这样几步。承担这套时序的是板上的 dwmmc 控制器，也就是新思科技
（DesignWare，一家提供可复用硬件电路设计的厂商）那套被广泛集成进各家芯片的 SD/eMMC 主机控
制器。一次数据传输完成后，控制器本应拉高一根中断线，通知内核这笔请求已经办妥；内核据此把
等待这笔请求的任务唤醒，去取回数据。这种由硬件主动报告完成、内核被动等待的方式称为中断驱
动，它让等待中的核心可以彻底休眠而不必空转轮询，是块设备完成路径的常规做法。Starry 的块运
行时正是按中断驱动来等待每一笔传输的完成。

把这块 SD 卡注册为块设备之后，启动却停在了挂载这一步。块运行时提交了第一笔 ext4 读取，随
即进入中断驱动的等待，可那根完成中断始终没有到来，等待便永远不会结束，挂载就此挂死。逐层
排查后确认问题出在控制器一侧。中断的路由与使能都核对无误，从设备树到中断控制器（GIC）的这
条线接得没错，控制器寄存器里的中断使能位也确实置上了，唯独这片 RK3588 上的 dwmmc 控制器在
一次数据传输结束后并不可靠地拉高它的完成中断。信号根本没有发出，内核这一侧再正确也无从被
唤醒，第一笔 ext4 读取于是永远不返回。

棘手之处在于，控制器是否真把中断送达，没有任何一个内核能直接读到的状态位能回答。运行时能
查询的 \texttt{is\_irq\_enabled()} 反映的只是控制器寄存器里那一位使能位，它说的是“我已被允
许发中断”，而非“中断确实经 GIC 送到了处理函数”。这两者在正常硬件上等价，在这片控制器上却
分了岔。既然没有直接的判据，唯一能观测到这一故障的地方，就是等待本身等了太久。

\textbf{据此在完成路径上补了一道对中断是否可用的判断。}中断驱动的等待不再无限期睡下去，而是带一个
有上限的期限。这条期限通过 \texttt{task\_wait\_until\_timeout} 实现，任务睡到完成条件成立即
被唤醒，或者睡满期限后自行醒来；它把底层带超时的等待原语的返回值取反，使得返回为真当且仅当
真正观测到了完成。期限取两秒，这是个有意放宽的值，普通慢速设备远不至于跑满，唯有完成信号确
实缺席才会触到。一旦在两秒内没有等到中断，运行时便判定这台设备此刻的完成中断不可用，先转为
轮询去亲自确认这笔传输究竟办妥没有，并把这台设备就地切到轮询，往后它的每一笔传输都直接读硬
件状态而不再等中断。这个切换是按设备一次性锁定的，由设备上一个记录完成方式的状态位
\texttt{completion\_mode} 经 \texttt{demote\_to\_polling} 翻成轮询，此后不再翻回，免得每一笔
传输都白白赔上两秒的期限。中断正常的设备完全不受影响，它们的完成中断照常即时唤醒等待者，期
限根本不会触发，也就没有任何额外开销。

\begin{lstlisting}[caption={完成等待的超时回退（节选自 \texttt{block\_runtime}）},captionpos=b]
// 中断驱动：带期限地睡，等完成或等满两秒
let completed = task_wait_until_timeout(
    || keys.iter().any(|&k| self.pending.lock().result(k).is_some()),
    IRQ_COMPLETION_FALLBACK_TIMEOUT, // 两秒
);
if !completed {
    // 期限内中断未到：判定本设备完成中断不可用，
    // 余生切到轮询，避免每笔传输都赔上两秒
    self.demote_to_polling();
    return self.polling_wait_for_any_key(&keys);
}
\end{lstlisting}

这道回退同时落在三条等待路径上，单笔请求的完成等待、队列推进的等待，以及面向多笔在途请求
的等待，任意一条等满两秒都会把设备切到轮询，缺席的中断因此不再能拖住挂载。设备出厂即以轮
询方式登记，只有在中断处理函数成功注册之后才被提升为中断驱动，这道回退正是这条提升的对称
反向，把一台被误判为中断可用的设备拨回它本就稳妥的轮询路径。

需要如实说明的是，这只是一道变通，不是根因修复。它保证的仅仅是挂载稳定可达，至于这片
dwmmc 控制器为何在传输结束后不拉高完成中断，仍是悬而未决的问题，根子在控制器一侧而非内核
的中断处理。真正的修复应当落在 dwmmc 主机驱动对控制器中断使能与中断屏蔽寄存器的设置上，
让控制器如约拉高那根线。在那之前，轮询回退让根文件系统能稳定挂上，整套系统得以启动，后续
各部件才有立足之地。\textbf{这项工作待提交至上游（提交 \texttt{5c36d90e3}）。}
